% sec:geo-null — Null Geodesics
\section{Null Geodesics}
\label{sec:geo-null}

Light rays follow the same extremal principle as massive particles, but with
a single additional constraint: the super-Hamiltonian vanishes identically
along the entire trajectory.  This constraint --- $\mathcal{H} = 0$ ---
is the defining equation of a null geodesic, and it is the reason the
Hamiltonian formulation is the natural choice for ray tracing.  The
Lagrangian system (\cref{eq:geo-lagrangian-first-order}) preserves the null
condition analytically but offers no cheap diagnostic when the integrator
violates it; the Hamiltonian system (\cref{eq:geo-xdot,eq:geo-pdot-cov})
turns any numerical violation of $\mathcal{H} = 0$ into a free accuracy
diagnostic.

\paragraph{The null constraint.}

A null geodesic is a curve $x^\mu(\lambda)$ whose tangent vector is
everywhere null: $\metric\,\dot{x}^\mu\dot{x}^\nu = 0$.  In Hamiltonian
language, this becomes
%
\begin{equation}
  \mathcal{H} = \tfrac{1}{2}\,\inversemetric\,p_\mu\,p_\nu = 0 \,.
  \label{eq:null-constraint}
\end{equation}
%
This is \cref{eq:geo-hamiltonian} evaluated on the null surface in
phase space.  Unlike the timelike case, where $\mathcal{H} = -\tfrac{1}{2}$
fixes the normalisation of the affine parameter to proper time, the null
constraint imposes no normalisation --- a freedom we exploit below.

\paragraph{Constraint preservation.}

Hamilton's equations preserve $\mathcal{H}$ exactly, without any special
effort.  To verify this, differentiate $\mathcal{H}$ along the geodesic
flow:
%
\begin{equation}
  \frac{\d\mathcal{H}}{\d\lambda}
    = \frac{\partial\mathcal{H}}{\partial x^\mu}\,\dot{x}^\mu
      + \frac{\partial\mathcal{H}}{\partial p_\mu}\,\dot{p}_\mu \,.
  \label{eq:H-dot}
\end{equation}
%
Substituting Hamilton's equations ---
$\dot{x}^\mu = \partial\mathcal{H}/\partial p_\mu$ and
$\dot{p}_\mu = -\partial\mathcal{H}/\partial x^\mu$ --- gives
%
\begin{equation}
  \frac{\d\mathcal{H}}{\d\lambda}
    = \frac{\partial\mathcal{H}}{\partial x^\mu}\,
      \frac{\partial\mathcal{H}}{\partial p_\mu}
      - \frac{\partial\mathcal{H}}{\partial p_\mu}\,
        \frac{\partial\mathcal{H}}{\partial x^\mu}
    = 0 \,.
  \label{eq:H-conservation}
\end{equation}
%
The two terms cancel identically.  This is not a special property of
the geodesic Hamiltonian or a consequence of any spacetime symmetry ---
it holds for any autonomous Hamiltonian system.  If $\mathcal{H} = 0$
at $\lambda = 0$, then $\mathcal{H} = 0$ for all $\lambda$.
\physnote{In the language of symplectic geometry, this is the statement
$\{\mathcal{H}, \mathcal{H}\} = 0$: every function Poisson-commutes
with itself.  The same argument shows that $\mathcal{H} = -\tfrac{1}{2}$
is preserved for timelike geodesics.}

The Lagrangian first-order system (\cref{eq:geo-lagrangian-first-order})
also preserves $\metric\,v^\mu v^\nu = 0$ analytically --- metric
compatibility guarantees this --- but provides no built-in scalar
diagnostic analogous to $\mathcal{H}$.  Under numerical discretisation,
both formulations suffer constraint drift; the Hamiltonian system makes
the drift immediately visible because $\mathcal{H}$ can be evaluated
from the state vector $(x^\mu, p_\mu)$ at negligible cost.  In the
Lagrangian system, the integrator provides no signal that the photon has
drifted off the null cone.
\warnnote{Symplectic integrators (St\"ormer--Verlet, leapfrog) preserve
$\mathcal{H}$ to within bounded oscillations rather than monotonic drift.
The codebase uses a non-symplectic Dormand--Prince integrator for its
higher order and adaptive step control, accepting slow $\mathcal{H}$ drift
in exchange for accuracy per step (\cref{sec:geo-conservation}).}

\paragraph{Affine rescaling freedom.}

Null geodesics have a parametrisation freedom that timelike geodesics lack.
For a massive particle, proper time $\tau$ provides a canonical choice of
affine parameter: $\metric\,\dot{x}^\mu\dot{x}^\nu = -1$ fixes the
normalisation.  For a photon, there is no rest frame and no proper time.
The reparametrisation $\lambda \to \alpha\lambda + \beta$ (with
$\alpha > 0$) sends one affine parameter to another without changing the
trajectory, and the covariant momentum scales as
%
\begin{equation}
  p_\mu \to \frac{1}{\alpha}\,p_\mu \,.
  \label{eq:affine-rescaling}
\end{equation}
%
Every conserved quantity built from $p_\mu$ scales the same way:
$E \to E/\alpha$, $L \to L/\alpha$.  But ratios of conserved quantities
are invariant, so the impact parameter $b = L/E$ and the reduced Carter
parameter $q = \sqrt{\mathcal{Q}}/E$ (defined for $\mathcal{Q} \geq 0$,
which holds for all null geodesics reaching a distant observer) are
independent of the parametrisation.  The trajectory itself --- which
spacetime points are visited, the deflection angle, whether the ray is
captured --- depends only on $(b, q)$.  This is why the ray tracer can
treat every pixel identically: the integration code needs only the
dimensionless ratios, not the energy scale.
\xrefnote{The $E = 1$ normalisation was introduced in
\cref{sec:sr-causal}; impact parameters $b$ and $q$ as observer-sky
coordinates were discussed in \cref{sec:geo-constants}.}

The standard choice is $E = 1$.  Given a photon direction
$(\theta, \varphi)$ on the observer's sky, the tetrad projection of
\cref{eq:null-tetrad} determines the initial null vector $k^\mu$ with
unit energy; lowering the index at the observer's position gives the
initial momentum $p_\mu = \metric\,k^\nu$.  Different pixels correspond
to different values of $(b, q)$, and the integration is identical for
each --- only the initial conditions differ.

\paragraph{Initial conditions from the null constraint.}

The tetrad construction automatically produces a null momentum, but it
is instructive to see how the constraint works in coordinates.

The Kerr--Schild inverse metric has a particularly simple structure that
makes the constraint equation explicit.  With
$\inversemetric = \eta^{\mu\nu} - f\,l^\mu l^\nu$
(\cref{eq:kerr-ks-inverse}), the null condition $\mathcal{H} = 0$ becomes
%
\begin{equation}
  \eta^{\mu\nu}\,p_\mu\,p_\nu - f\,(l^\mu p_\mu)^2 = 0 \,.
  \label{eq:null-constraint-ks}
\end{equation}
%
Given three components of $p_\mu$ (from the pixel direction and the
$E = 1$ normalisation), this is generically a quadratic in the remaining
component, yielding two solutions.  Since the ray tracer works backward
--- tracing from camera to source --- the physical root is the one that
sends the ray into the past along the observer's worldline.  Once $p_\mu$
is set, the initial state $(x^\mu_0, p_{\mu,0}, \lambda_0 = 0)$ is
complete, and the integrator takes over.
\implnote{The \texttt{HamiltonianState} record holds $(x^\mu, p_\mu,
\lambda)$ as the full state vector.  The function \texttt{hamiltonian}
evaluates $\mathcal{H}$ from the inverse metric and current momentum,
providing the constraint diagnostic at every accepted step.}

\paragraph{The codebase implementation.}
\coderef{Geodesic}{Hamiltonian}

The function \texttt{hamiltonianRHS} in
\codemodule{Geodesic.Hamiltonian} computes the full eight-component
right-hand side $(\dot{x}^\mu, \dot{p}_\mu)$ of the geodesic ODE system.
Its type signature exposes the design:
%
\begin{lstlisting}[language=Haskell]
hamiltonianRHS
  :: (forall a. Floating a => [a] -> [a])
  -> Sym4 'Contravariant
  -> HamiltonianState
  -> (Vec4 'Contravariant, Vec4 'Covariant)
\end{lstlisting}
%
The rank-2 polymorphic first argument is the metric function --- passing
it an AD type instead of \texttt{Double} produces exact metric
derivatives, a technique developed in \cref{ch:autodiff}.  The inverse
metric (\texttt{Sym4 'Contravariant}) is precomputed once per step and
passed in, since the $4 \times 4$ inversion is cheap but wasteful to
repeat.  The return type is the pair $(\dot{x}^\mu, \dot{p}_\mu)$ ---
contravariant velocity and covariant force --- ready for the adaptive
integrator.

A companion function \texttt{hamiltonian} evaluates the
super-Hamiltonian $\mathcal{H}$ from the inverse metric and current
momentum.  The codebase calls this after
every accepted integration step; for a typical ray trace through the Kerr
geometry with tolerances of $10^{-10}$, the violation remains below
$\abs{\mathcal{H}} < 10^{-12}$ throughout the trajectory.  The null
constraint, which both formulations preserve analytically, becomes a
continuous numerical accuracy certificate --- requiring no extra ODE
solves, just a single scalar evaluation per step.
